% Setting Up, Chapter 1 -- Before You Begin (annex Sec 1, SU: check the
% package, the hardware, the box's Hardware Chart [lot s1; sur s2]).
% Sources (Law 1): README.md System requirements; spec/hardware.json (cpu,
% memory, fpu.required_by ["Initech 123"]); ADR-0009 Amendment DEC-01a
% AM-3 (Initech 123 requires a 486DX or a 387/487); the shipped data
% volume (the Makefile's flair-data-rules: APPS/123.EXE and APPS/123.MNU).
% GATES: test-flair-i123 test-flair-i123-bochs test-pubs-build
\Gates{test-flair-i123, test-flair-i123-bochs, test-pubs-build}
\chapter{Before You Begin}

This chapter tells you what you need before you use Initech 123: the pieces
of the package, the computer and its equipment, and a chart in which to
record your equipment.

\section{What Comes with Initech 123}
\ix{package}
Initech 123 is supplied with InitechOS.  Its program is installed with the
operating system: the files \file{123.EXE} (the program) and \file{123.MNU}
(its menus) are in the folder \file{APPS} of the InitechOS data disk.  You do
not install Initech 123 separately.

The Initech 123 documentation consists of:
\begin{itemize}
\item this volume, \doctitle{Setting Up Initech 123, Tutorial, Quick Start
  and Sample Applications};
\item the \doctitle{Initech 123 Reference};
\item the \doctitle{Initech 123 Quick Reference};
\item two keyboard templates, one for the 101-key and one for the 84-key
  keyboard.
\end{itemize}
\draft{PUB-REF-0001 PK-08..PK-13: the box, the license envelope and the
cards are planned and not yet produced}

\section{What Your Computer Needs}
\ix{hardware requirements}
Initech 123 runs on any computer that runs InitechOS and has a numeric
coprocessor.
\begin{itemize}
\item An Intel 80486DX processor, which has the numeric coprocessor built
  in; or an 80386 processor with an 80387 coprocessor, or an 80486SX with an
  80487.
\item The memory InitechOS requires: at least 4 MB of extended memory above
  the first megabyte.
\item A VGA display and monitor.
\item A fixed disk and a diskette drive, as InitechOS requires.
\item A mouse, for starting Initech 123 from the FLAIR desktop.
\end{itemize}

\begin{Note}
Initech 123 does all its arithmetic with the numeric coprocessor.  It does
not run on a computer without one, even though InitechOS itself does.
\end{Note}

\section{The Hardware Chart}
\ix{Hardware Chart}
Record your equipment in this chart.  Keep it with your Initech 123 books:
your Initech representative will ask for it if you call for help.

\par\addvspace{6pt}
\noindent\begin{tabular}{@{}p{58mm}p{70mm}@{}}
\hline
\sffamily\small\textbf{Item} & \sffamily\small\textbf{Your equipment}\\
\hline
\rule{0pt}{14pt}Computer, make and model & \\ \hline
\rule{0pt}{14pt}Processor (80486DX, 80386 and 80387, ...) & \\ \hline
\rule{0pt}{14pt}Memory (megabytes) & \\ \hline
\rule{0pt}{14pt}Display adapter and monitor & \\ \hline
\rule{0pt}{14pt}Fixed disk (megabytes) & \\ \hline
\rule{0pt}{14pt}Diskette drives (size and capacity) & \\ \hline
\rule{0pt}{14pt}Mouse & \\ \hline
\rule{0pt}{14pt}Printer & \\ \hline
\rule{0pt}{14pt}InitechOS serial number & \\ \hline
\rule{0pt}{14pt}Date of purchase & \\ \hline
\end{tabular}
